\section{Applicable standards and how they are identified}
To identify the applicable harmonised EMC emission standards, the European Commission~[\ref{ref:ec-emc-directive}]
webpage for harmonised standards under the EMC Directive 2014/30/EU was used as the
starting point. This page provides a ``Summary list of titles and references of harmonised
standards under Directive 2014/30/EU for EMC'', available as a downloadable PDF or
spreadsheet. The summary list contains, among other information, the reference number of
each standard and the title of the standard.

The list was first reviewed by comparing the standard titles with the intended environment
and function of the AGV.

The following candidate standards were selected from the summary list for further
comparison:

\begin{table}[H]
\centering
\caption{Candidate harmonised EMC standards selected from the summary list under Directive 2014/30/EU.}
\small
\setlength{\tabcolsep}{4pt}
\begin{tabular}{@{}p{0.32\linewidth}p{0.62\linewidth}@{}}
\hline
{\raggedright\textbf{Standard}\par} & {\raggedright\textbf{Title in the summary list}\par} \\
\hline
{\raggedright EN~61800-3:2004+A1:2012\par} &
{\raggedright Adjustable speed electrical power drive systems, Part 3: EMC requirements and specific test methods\par} \\

{\raggedright EN~61000-6-3:2007+A1:2011\par} &
{\raggedright Electromagnetic compatibility (EMC), Part 6-3: Generic standards, Emission standard for residential, commercial and light-industrial environments\par} \\

{\raggedright EN~61000-6-4:2007+A1:2011\par} &
{\raggedright Electromagnetic compatibility (EMC) - Part 6-4: Generic standards - Emission standard for industrial environments\par} \\

{\raggedright EN~14010:2003+A1:2009\par} &
{\raggedright Safety of machinery, Equipment for power driven parking of motor vehicles, Safety and EMC requirements for design, manufacturing, erection and commissioning stages\par} \\
\hline
\end{tabular}
\end{table}

The best fit is EN~14010:2003+A1:2009, \textit{Safety of machinery --- Equipment for power driven parking of motor vehicles --- Safety and EMC requirements for design, manufacturing, erection and commissioning stages}.

As this standard is licence-based, it cannot be referenced in detail directly in this
report. However, EN~14010 does not itself define a separate dedicated EMC emission
limit set. Instead, it states that the electromagnetic disturbances generated by the power
driven parking equipment shall not exceed the levels specified in the generic emission
standard EN~61000-6-3.

For this reason, the practical EMC focus in this report is placed on
EN~61000-6-3, \textit{Electromagnetic compatibility (EMC), Part 6-3: Generic standards,
Emission standard for residential, commercial and light-industrial environments}.

Since EN~61000-6-3 is also licence-based, this test report does not claim formal
conformity with the standard. The standard is therefore not used directly as a complete
test method document in this report. Instead, it is used only as a practical
pre-compliance reference, and the numerical limits used in the measurements are taken
from publicly available secondary sources that reference EN~61000-6-3.

\begin{table}[H]
\centering
\caption{Secondary sources used for practical radiated-emission reference limits.}
\small
\setlength{\tabcolsep}{4pt}
\begin{tabular}{@{}p{0.25\linewidth}p{0.24\linewidth}p{0.16\linewidth}p{0.27\linewidth}@{}}
\hline
{\raggedright\textbf{Frequency range}\par} &
{\raggedright\textbf{Limit}\par} &
{\raggedright\textbf{Page}\par} &
{\raggedright\textbf{Source}\par} \\
\hline

{\raggedright \(30\,MHz\) to \(230\,MHz\)\par} &
{\raggedright \(40\,dB\mu V/m\) QP at \(3\,m\)\par} &
{\raggedright p.~8 of 20\par} &
{\raggedright BK Services EMC report~[\ref{ref:bk-superchrono}]\par} \\

{\raggedright \(230\,MHz\) to \(1000\,MHz\)\par} &
{\raggedright \(47\,dB\mu V/m\) QP at \(3\,m\)\par} &
{\raggedright p.~8 of 20\par} &
{\raggedright BK Services EMC report~[\ref{ref:bk-superchrono}]\par} \\

{\raggedright \(30\,MHz\) to \(230\,MHz\)\par} &
{\raggedright \(40\,dB\mu V/m\) QP at \(3\,m\)\par} &
{\raggedright p.~33\par} &
{\raggedright Vecow, CE EMC Test Report~[\ref{ref:vecow-ce-emc}]\par} \\

{\raggedright \(230\,MHz\) to \(1000\,MHz\)\par} &
{\raggedright \(47\,dB\mu V/m\) QP at \(3\,m\)\par} &
{\raggedright p.~33\par} &
{\raggedright Vecow, CE EMC Test Report~[\ref{ref:vecow-ce-emc}]\par} \\

{\raggedright \(30\,MHz\) to \(230\,MHz\)\par} &
{\raggedright \(40\,dB\mu V/m\) QP at \(3\,m\)\par} &
{\raggedright p.~8 of 17, section~6.1\par} &
{\raggedright Adeo, EMC Test Report, EN~61000-6-3~[\ref{ref:adeo-emc}]\par} \\

{\raggedright \(230\,MHz\) to \(1000\,MHz\)\par} &
{\raggedright \(47\,dB\mu V/m\) QP at \(3\,m\)\par} &
{\raggedright p.~8 of 17, section~6.1\par} &
{\raggedright Adeo, EMC Test Report, EN~61000-6-3~[\ref{ref:adeo-emc}]\par} \\

\hline
\end{tabular}
\end{table}

Based on the secondary sources, it can be inferred that the practical
EN~61000-6-3 radiated-emission reference limits at \(3\,m\) are
\(40\,dB\mu V/m\) from \(30\,MHz\) to \(230\,MHz\), and
\(47\,dB\mu V/m\) from \(230\,MHz\) to \(1000\,MHz\), using a quasi-peak
detector.



